\documentclass[10pt, twocolumn]{article}

% A two-column paper in a generic conference style. If your venue gives you
% its own class or style file, add it to the project and switch to it.

% Fonts and encoding: pdfLaTeX needs these; XeLaTeX and Tectonic do not.
\usepackage{iftex}
\ifPDFTeX
  \usepackage[T1]{fontenc}
  \usepackage{lmodern}
\fi
\usepackage{microtype}
\usepackage{amsmath, amssymb}
\usepackage{graphicx}
\usepackage{booktabs}
\usepackage[letterpaper, margin=0.75in, columnsep=0.3in]{geometry}
\usepackage{xcolor}
\usepackage[colorlinks=true, linkcolor=blue!60!black, citecolor=green!40!black, urlcolor=blue!60!black]{hyperref}

\title{\Large\bfseries A Concise and Informative Title for a Conference Paper}
\author{
  First Author \\
  \small Department, University \\
  \small City, Country \\
  \small \texttt{first@example.org}
  \and
  Second Author \\
  \small Research Laboratory \\
  \small City, Country \\
  \small \texttt{second@example.org}
}
\date{}

\begin{document}

\maketitle

\begin{abstract}
State the problem, your approach and your main result in about 150 words.
Many reviewers read only the title and the abstract, so say plainly what
is new and how much better it is than what came before.
\end{abstract}

\noindent\textbf{Keywords:} first keyword, second keyword, third keyword

\section{Introduction}
\label{sec:intro}

Explain the problem and why the community should care. Summarise the
closest earlier work~\cite{surname2023conf, other2022journal} and what it
leaves open. Most venues expect a clear list of contributions:
\begin{itemize}
  \item the first contribution, stated as a result rather than an activity;
  \item the second contribution;
  \item an evaluation on a public benchmark (Section~\ref{sec:eval}).
\end{itemize}

\section{Related Work}
\label{sec:related}

Group earlier work by approach and say how yours differs from each group.
Cite the strongest competing method in each group~\cite{lamport1994}
and state where it falls short.

\section{Approach}
\label{sec:approach}

Define the notation before you use it. Given inputs
$x_1, \dots, x_n$ and weights $w_i$, the model predicts
\begin{equation}
  \hat{y} = \sigma\Bigl( \sum_{i=1}^{n} w_i x_i + b \Bigr) ,
  \label{eq:model}
\end{equation}
where $\sigma$ is the logistic function. The weights minimise the
average loss over $N$ training examples,
\begin{equation}
  L(w) = \frac{1}{N} \sum_{j=1}^{N} \ell(y_j, \hat{y}_j) ,
  \label{eq:loss}
\end{equation}
where $\ell$ is the cross-entropy between a label and its prediction.

\begin{figure*}[t]
  \centering
  \fbox{\parbox[c][3.5cm][c]{0.8\textwidth}{\centering Wide figure placeholder: an overview of the approach}}
  \caption{A figure that spans both columns. Use \texttt{figure*} for wide diagrams and plain \texttt{figure} for ones that fit a single column.}
  \label{fig:overview}
\end{figure*}

\section{Evaluation}
\label{sec:eval}

Describe the data, the baselines and the metrics, then the results.
Figure~\ref{fig:overview} gives an overview of the approach in
Section~\ref{sec:approach}, and Table~\ref{tab:results} compares it with
two baselines using Equation~\eqref{eq:loss} as the training objective.

\begin{table}[t]
  \centering
  \small
  \begin{tabular}{lrr}
    \toprule
    Method & Accuracy (\%) & Time (s) \\
    \midrule
    Baseline A & 81.2 & 14 \\
    Baseline B & 84.5 & 32 \\
    Ours & \textbf{88.9} & 17 \\
    \bottomrule
  \end{tabular}
  \caption{Results on the test set (placeholder numbers).}
  \label{tab:results}
\end{table}

\section{Conclusion}
\label{sec:conclusion}

Summarise what the paper showed, its limitations, and the most promising
direction for future work.

\section*{Acknowledgements}

Thank funders and colleagues here. Leave this section out of an anonymous
submission.

\bibliographystyle{plain}
\bibliography{references}

\end{document}
